\chapter{Literature formulas as limits}\label{app:F}
Each row of \cref{tab:F-rows} records a published formula that the book recovers as a limit, together with four items that its Limit Check does not state: the source equation, the approximation chain, the conditions of validity and the leading correction. The first column names the formula, its equation in the book and its Limit Check, or marks the row ``(no box)''; \cref{sec:9.4} explains how a box is read and places the formulas of \crefrange{chap:10}{chap:17} on the truncation label. The rows follow the order of the chapters, and the last column gives each formula's home section, where it is derived and where its box, if it has one, stands. The Leading-correction column opens with the size of that correction in bold, taken from the text, the Limit Check or its figure; \Fsize{unsized} marks a correction that the book does not yet size.

% One row per literature formula recovered as a limit: source equation, size of the leading correction, home section.
{\footnotesize
\setlength{\tabcolsep}{4pt}
\renewcommand{\arraystretch}{1.1}
\begin{longtable}{@{}>{\raggedright\arraybackslash}p{0.36\textwidth}>{\raggedright\arraybackslash}p{0.19\textwidth}>{\raggedright\arraybackslash}p{0.30\textwidth}>{\raggedright\arraybackslash}p{0.07\textwidth}@{}}
\caption{Literature formulas recovered as limits: source equation, size of the leading correction and home section.}\label{tab:F-rows}\\
\toprule
Formula & Source (eq.) & Leading correction & Section \\
\midrule
\endfirsthead
\toprule
Formula & Source (eq.) & Leading correction & Section \\
\midrule
\endhead
\bottomrule
\endlastfoot
Elastic nuclear recoil (\cref{lim:1.3a}), \cref{eq:1.3-ERmax,eq:1.3-vmin}
& \cite[Eqs.~(77), (78)]{lin2019tasi}
& relative width $(\omega_0/E_R)^{1/2} \le 1/3$ for $E_R \ge \SI{0.54}{eV}$ (Si)
& \cref{sec:1.3} \\
Heavy and light mediators, $\FDM = 1$ and $(q_0/q)^2$ (\cref{lim:1.4a}), \cref{eq:1.4-yukawa-sigma}
& \cite[Eq.~(15) and the display above it]{essig2012direct}
& $\le$\,\SI{10}{\percent} for $m_V \ge 4.5\,q_0$ (heavy) or $m_V \le 0.026\,q_0$ (light)
& \cref{sec:1.4} \\
Atomic ionization rate (\cref{lim:2.2a}), \cref{eq:2.2-ion,eq:2.2-fion}
& \cite[Eqs.~(5), (6); Eq.~(1) for the Fermi factor]{essig2012direct}
& about \SI{30}{\percent} (outer shells of heavy atoms); $\times0.46$ (H 1s, \SI{1}{GeV})
& \cref{sec:2.2} \\
Thomas--Fermi screening (\cref{lim:4.5a}), \cref{eq:4.4-TF}
& \cite[Eq.~(S.14)]{hochberg2021determining}
& $\le$\,\SI{2}{\percent} for $q < 0.49\,k_{\mathrm{F}}$; \SI{10}{\percent} at $1.06\,k_{\mathrm{F}}$
& \cref{sec:4.5} \\
Plasmon dispersion and cutoff (\cref{lim:4.5b}), \cref{eq:4.5-dispersion}
& \cite[Eq.~(171) and the text after it]{kahn2022searches}
& $\le$\,\SI{1}{\percent} for $q < 0.39\,k_{\mathrm{F}}$ ($r_s = 2$); cutoff estimate \SI{9}{\percent} low
& \cref{sec:4.5} \\
Gapless Dirac dielectric function (\cref{lim:7.3a}), \cref{eq:D.22-gapless}
& \cite[Eq.~(2.7)]{hochberg2018detection}, whose cutoff is a hard cutoff; $g_{\mathrm{deg}}$ taken out of the coupling
& $\le 0.1$ (units of $g_{\mathrm{deg}}\tilde\alpha/6\pi$) for $v_{\mathrm{D}}q > 7.5\,\Delta_{\mathrm{D}}$
& \cref{sec:7.3} \\
Second Born correction for a Yukawa potential (\cref{lim:8.2a}), \cref{eq:8.2-born}
& \cite[Chap.~6]{sakurai2017modern} (conditions $|f|m_V \ll 1$ and $y_\chi y_T/$ $4\pi v \ll 1$, \cref{sec:2.2})
& $\le$\,\SI{30}{\percent} of $\delta_{\mathrm{Born}}$ for $p_{\mathrm{rel}} \le m_V$; $|\delta_{\mathrm{Born}}| \le 8.0\times10^{-5}$ at the benchmark cross-sections
& \cref{sec:8.2} \\
Crystal form-factor rate (\cref{lim:10.1a}), \cref{eq:10.1-fcrys}
& \cite[Eqs.~(3.13) and (3.17)]{essig2016direct}, divided by the target mass $N_{\mathrm{cell}}M_{\mathrm{cell}}$
& screening $\times0.47$ ($q$-integrated, \SI{5}{eV}); 0.05--0.61 pointwise
& \cref{sec:10.1} \\
Loss-function rate (\cref{lim:10.2a}), \cref{eq:10.2-rate}
& \cite[Eqs.~(6) and (13)]{knapen2021dark}; \cite[Eq.~(1)]{hochberg2021determining}
& inside the phonon band $>10^4\times$ the electronic rate (electron-only coupling, \SI{1}{MeV}, model)
& \cref{sec:10.2} \\
Sum-rule bounds: $f$-sum and screening sum rules, rate bound (\cref{lim:10.2b}), \cref{eq:10.2-bound}
& \cite[Eqs.~(1), (8) and (9)]{lasenby2022dark}; \cite[Eqs.~(1), (2), (8) and (9)]{barker2026first}; non-density couplings \cite[Eqs.~(4.16)--(4.20)]{catena2025general}
& relaxed-lattice term $\mathcal{G}_{\mathrm{rel}} = 1$ at \SI{18}{keV} (Si), growing as $q^{-2}$ below
& \cref{sec:10.2} \\
Electron-gas (Lindhard, Mermin) loss-function rate (\cref{lim:10.3a}), \cref{eq:10.3-gasrate,eq:4.5-lindhard}
& \cite[Eqs.~(13), (17)--(19); oscillator form Eqs.~(20), (21)]{knapen2021dark}; \cite{knapen2022darkelf}
& Lindhard/silicon 0.70--1.48 at 4.2--\SI{9.5}{eV}; $\times24$ at \SI{2}{eV}
& \cref{sec:10.3} \\
Plasmon production from the pole of $1/\varepsilon$, light mediator (\cref{lim:10.3b}), \cref{eq:10.3-plrate}
& \cite[Eqs.~(5), (7)]{kozaczuk2020plasmon}; \cite[Eqs.~(4), (7), (13), (14)]{kurinsky2020dark}
& pole share 0.99 at $q_{\mathrm{c}}/2$, 0.77 at $0.9\,q_{\mathrm{c}}$; cutoff choice $\times1.24$
& \cref{sec:10.3} \\
Velocity coupling $\mathcal{O}_8$ with the electron at rest; crystal responses $\overline W_1$--$\overline W_5$ (\cref{lim:10.4a}), \cref{eq:10.4-O8}
& \cite[Eqs.~(5), (6)]{hochberg2025determining}; \cite[Eqs.~(36)--(42)]{catena2021crystal}; \cite[Eqs.~(18), (23), App.~B]{trickle2023exceed}; \cite[Table~III]{trickle2022effective}
& dropped transverse term $r_{\mathrm{IP}} = 22$--56 times the kept one (Si, 3--\SI{30}{MeV})
& \cref{sec:10.4} \\
Bethe--Salpeter exciton rate (\cref{lim:10.5a}), \cref{eq:10.5-rate}
& \cite[Eqs.~(1) and (2)]{taufertshofer2026excitonic}
& about \SI{10}{\percent} on $\varepsilon_2$ (coupling block)
& \cref{sec:10.5} \\
Absorption with in-medium mixing: dark photon; axion-like particle in a field (\cref{lim:10.7a}), \cref{eq:10.7-dp,eq:10.7-alpB}
& \cites[Eqs.~(6) and (7)]{hochberg2017absorption}[Eqs.~(2), (18) and (20)]{an2015direct}[Eqs.~(38)--(40)]{mitridate2021dark}[Eqs.~(3), (6) and (7)]{berlin2024absorption}
& $O(v^2)$: $-1.7\times10^{-6}$ at \SI{10}{eV}; vacuum mixing off by $4\times10^{-4}$ to 20 below \SI{63}{eV}
& \cref{sec:10.7} \\
Scalar absorption through the kinetic-energy row and its sum rule (\cref{lim:10.7b}), \cref{eq:10.7-scalar,eq:10.7-sumrule}
& \cite[Eqs.~(18), (55) and (56)]{mitridate2021dark}, with $\bar v^2 = \mathcal{T}_e/m_e$ and $g_{\phi ee} =$ $d_{\phi ee}\sqrt{4\pi}$ $\times m_e/M_{\mathrm{Pl}}$
& density row $3\times10^{-3}$ (\SI{5}{eV}) to about 1 (16--\SI{18}{eV}) of the kinetic rate
& \cref{sec:10.7} \\
Pair absorption of a complex scalar (\cref{lim:10.8a}), \cref{eq:10.8-contact}
& \cite[Eqs.~(3) and (5)--(8)]{bauer2026dark}; $\xi = 44.8$ gives $\rho_\chi/2m_\phi$ per species
& $\mathrm{Var}(K)/\gamma^2$: $\le$\,\SI{10}{\percent} for $\gamma \ge 1.5\,\bar K$
& \cref{sec:10.8} \\
Dark Compton scattering of a vector on a free electron at rest (\cref{lim:10.8b}), \cref{eq:10.8-compton,eq:10.8-thomson}
& \cite[Eq.~(2) and App.\ Eqs.~(7)--(12)]{hochberg2022impact}, their $\varepsilon$ = our $\zeta$
& $\tfrac58x^3$: $\le$\,\SI{1}{\percent} for $m_\phi < \SI{165}{keV}$, \SI{10}{\percent} below \SI{517}{keV}
& \cref{sec:10.8} \\
Dark Compton scattering of an axion-like particle (\cref{lim:10.8b}): published rate, against \cref{eq:10.8-alpC}
& \cite[Eqs.~(3) and (17)]{hochberg2022impact}; squared amplitude \cite[Eq.~(15)]{hochberg2022impact}
& published form high by $\times1.836$ at \SI{100}{keV} ($\times2$ as $x \to 0$); spin seagull relative $O(x)$
& \cref{sec:10.8} \\
Inelastic electron recoil in crystals (\cref{lim:10.9a}), \cref{eq:10.9-inel,eq:10.9-shell}
& \cite[Eqs.~(10), (11), (15), (16) and (20)]{gu2022detection}, with the factor $q^{-2}$ of \cite[Eq.~(3.13)]{essig2016direct} restored
& screening $\times9.4$--18.5 at \SI{1}{MeV} (heavy), up to 317 (light); shell form $\approx 2.5\langle v^2\rangle/v_\Delta^2$
& \cref{sec:10.9} \\
Boosted dark matter, relativistic density row (\cref{lim:10.9b}), \cref{eq:10.9-Drel,eq:10.9-bound}
& \cites[Eqs.~(4)--(6)]{liang2025plasmon}[Eqs.~(9)--(11)]{aggarwal2026constraints}[Eq.~(2)]{liang2026dielectric}
& transverse row $2.7\times10^{-5}$ ($v = 0.1$) to \SI{5.5}{\percent} ($v = 0.99$) of the silicon rate
& \cref{sec:10.9} \\
Fermionic absorption on electrons (\cref{lim:10.9c}), \cref{eq:10.9-ferm}
& \cites[Eqs.~(7)--(10) and (30)]{dror2021absorption}[Eqs.~(A.18)--(A.21)]{essig2016direct}[Eqs.~(1), (2), (14) and (15)]{cox2024absorption}
& unsized
& \cref{sec:10.9} \\
In-medium decay of an excited state, $\tau^{-1} \sim \pi|\Delta_\chi|^3/\Lambda_d^2$ (\cref{lim:10.9d}), \cref{eq:10.9-vac,eq:10.9-decay}
& \cite[Eqs.~(2.1) and (2.2)]{baryakhtar2020electromagnetic}
& unsized
& \cref{sec:10.9} \\
Long-wavelength acoustic structure factor (\cref{lim:11.1a}), \cref{eq:11.1-Sac}
& \cite[Eq.~(30) with Eq.~(27)]{campbelldeem2022dark}, longitudinal term, with $m_p$ for $m_u$ and the cutoff $\Theta(\omega_{\mathrm{LO}} - \omega)$; prior art \cite[Eq.~(10)]{campbelldeem2020multiphonon}
& $\le$\,\SI{10}{\percent} for $q < 0.48\,q_{\mathrm{BZ}}$; \SI{26}{\percent} at $q_{\mathrm{BZ}}$
& \cref{sec:11.1} \\
Light-mediator single-phonon rate per mass (\cref{lim:11.1b}), \cref{eq:11.1-light}
& \cite[Eq.~(35)]{griffin2020multichannel}, with their $\bar\sigma_n$ at the reference momentum $\mchi v_0$
& $(1 + \kappa_\eta)m_{\mathrm{min}}^{\mathrm{ac}}/\mchi$: $\le$\,\SI{10}{\percent} for $\mchi \ge 44\,m_{\mathrm{min}}^{\mathrm{ac}}$
& \cref{sec:11.1} \\
Sphere-averaged longitudinal speed $\bar c_{\mathrm{LA}}$ (\cref{lim:11.1c}), \cref{eq:11.1-cbarLA}
& \cite[App.~A\,3]{griffin2020multichannel} (there from density-functional $\omega_{\mathrm{LA}}/q$ on a sphere of radius \SI{10}{eV})
& $\le$\,\SI{1}{\percent} (Si 0.67, Ge 0.92, diamond \SI{0.21}{\percent})
& \cref{sec:11.1} \\
Born-charge one-phonon structure factor of a total-charge coupling (\cref{lim:11.2a}), \cref{eq:11.2-SZZ}
& \cites[Eqs.~(92), (105), (106)]{trickle2020multichannel}[Eqs.~(19), (23), v1 numbering]{griffin2018directional}[Eq.~(61)]{trickle2022effective}; one branch: \cite[Eqs.~(4)--(7)]{knapen2018detection}, whose $C_F^2 \propto 1/V_{\mathrm{cell}}$ must be read with the normalization volume, which cancels
& $O(qa)$: \SI{1}{\percent} at \SI{90}{eV}, \SI{5}{\percent} at \SI{450}{eV} (GaAs, $[111]$)
& \cref{sec:11.2} \\
One polar branch and the band-integral rate (\cref{lim:11.2b}), \cref{eq:11.2-Gamma,eq:11.2-Ipol,eq:11.2-diatomic}
& \cite[Eqs.~(21), (22), (24), v1 numbering]{griffin2018directional} (prior art for the one-branch formula; the band-integral form is \cref{eq:11.2-Ipol})
& $\le$\,\SI{7}{\percent} above \SI{64}{keV} (sapphire); $\times2.1$ at \SI{40}{keV}
& \cref{sec:11.2} \\
Static dielectric constant of one polar branch (\cref{lim:11.2c}), \cref{eq:11.2-static}
& \cites[Chap.~II]{born1954dynamical}[Eq.~(B14) at $\omega = 0$, v1 numbering]{griffin2018directional}
& exact for cubic diatomics; in sapphire the rate branch holds \SI{2.7}{\percent} of $\varepsilon_0 - \varepsilon_\infty$
& \cref{sec:11.2} \\
Born charge over $\varepsilon_\infty$ from the LO--TO splitting (\cref{lim:11.2d}), \cref{eq:11.2-Zeps}
& \cites[Chap.~II]{born1954dynamical}[Eqs.~(B11), (B15), v1 numbering]{griffin2018directional}; \cite{lyddane1941polar}
& \SI{1.0}{\percent} for \SI{0.1}{\percent} frequency and \SI{1}{\percent} $\varepsilon_\infty$ errors (GaAs)
& \cref{sec:11.2} \\
Multiphonon incoherent approximation and Poisson tail (\cref{lim:11.3a}), \cref{eq:11.3-inc}
& \cites[Eqs.~(15), (23), (24), (36)]{campbelldeem2022dark}[Eqs.~(7), (10)]{kahn2021dark}[Eq.~(12)]{campbelldeem2020multiphonon}
& $\delta^{\mathrm{coh}}_2 = +0.35$ ($q \to 0$), $-0.11$ ($q_{\mathrm{BZ}}$), $-0.018$ (\SI{10}{keV})
& \cref{sec:11.3} \\
Second-order thermal diffuse scattering (\cref{lim:11.3b}), \cref{eq:11.3-tds}
& \cite{warren1969xray} (classical $1/|\vect{q} - \vect{G}|$ law)
& classical coefficient 0.46 of the bound; $\le$\,\SI{10}{\percent} for $|\vect{q} - \vect{G}| \le 0.13\,q_{\mathrm{BZ}}$ (\SI{293}{K})
& \cref{sec:11.3} \\
Free nuclear recoil with the Helm form factor (\cref{lim:11.4a}), \cref{eq:11.4-LS,eq:2.2-helm}
& \cites[Eqs.~(37)--(39), (42)]{campbelldeem2022dark}[Eq.~(104)]{lin2019tasi}{lewin1996review}; \cites[Eqs.~(90)--(92)]{kahn2022searches}[Eqs.~(4), (12)]{kahn2021dark}
& relative width $x_R^{-1/2}$: $\le$\,\SI{10}{\percent} for $x_R \ge 2.3$, \SI{2}{\percent} for $x_R \ge 4.3$ (Si, \SI{1}{GeV})
& \cref{sec:11.4} \\
Fermionic absorption on nuclei (\cref{lim:11.4b}), \cref{eq:11.4-ferm-nuc}
& \cite[Eqs.~(1), (3), (4)]{dror2020direct} (their $\sigma_{\mathrm{NC}}$, $\rho_\chi \simeq \SI{0.4}{GeV.cm^{-3}}$)
& $\le$\,\SI{10}{\percent} for \SI{61}{keV}--\SI{3.0}{GeV}; zero-phonon weight 0.94 at \SI{10}{keV}
& \cref{sec:11.4} \\
Anharmonic corrections to multiphonon weights (\cref{lim:11.5a}), \cref{eq:11.5-toy,eq:11.5-two,eq:11.5-anh}
& \cite[Eqs.~(26), (39), (45)--(48), (65)]{lin2024anharmonic}; \cite[Eq.~(13)]{campbelldeem2020multiphonon}
& $\le$\,\SI{10}{\percent} for $q \ge \SI{2}{keV}$ (Si model); impulse width $-0.14$ to $+\SI{0.08}{\percent}$
& \cref{sec:11.5} \\
Solid helium, single-phonon rate, quasi-harmonic and self-consistent (\cref{lim:11.5b}), \cref{eq:11.1-S1,eq:5.7-scha}
& \cite[Eqs.~(2), (4), (5)]{ashour2024pressure}; \cite[Table~I]{cazorla2015first} (reference Lindemann ratios)
& Gaussian (SCHA) miss $\le$\,\SI{10}{\percent} for $r_{\mathrm{L}} \le 0.227$, 0.896 of exact at $r_{\mathrm{L}} = 0.230$ (cell model)
& \cref{sec:11.5} \\
Markov--Krein band and its tail (\cref{lim:11.5c}), \cref{eq:C.10-band,eq:C.10-rho,eq:C.10-tail}
& \cites{krein1977markov}{akhiezer1965classical} (book level; classical moment theory)
& tail form exceeded by \SI{34}{\percent} at $\xi = 3$, \SI{20}{\percent} at $\xi = 6$
& \cref{sec:11.5} \\
Incoherent spin absorption (\cref{lim:11.6a}), \cref{eq:11.6-alp}
& \cite[Eqs.~(25)--(31)]{bloch2025broadband} (their $g_N = g_{aNN}/2m_N$, $g_p\lambda_j = g_{ai}/2m_N$, $\xi_j$ the isotope-averaged $\lambda_j^2J(J+1)/M$, $D_j = g_{\mathrm{ph},\kappa}/3$)
& $-3.5$ to $+\SI{10}{\percent}$ below \SI{55}{meV} (Si); 1.78 and 0.72 near the top of the optical band
& \cref{sec:11.6} \\
Lifshitz mass-defect Green function (\cref{lim:11.6b}), \cref{eq:11.6-lifshitz}
& \cites{lifshitz1956some}{dawber1963vibrations} (classical result; equations not verified against the original texts)
& pair term $\le$\,\SI{2.2}{\percent} (natural Si, to \SI{55}{meV}); $g_{\mathrm{md}} = 2.0\times10^{-4}$
& \cref{sec:11.6} \\
Daily modulation from the anisotropic response (\cref{lim:11.7a}), \cref{eq:11.7-rate,eq:11.7-ups}
& \cites[Eqs.~(6), (10), (17), (25), (37)--(41)]{aladulrazzaq2026anisotropic}[Eqs.~(18)--(24), (28)--(31)]{stratman2025daily}
& $\mathcal{H}_6/\mathcal{H}_0 \lesssim 10^{-2}$ for $\vmin \lesssim \SI{600}{km.s^{-1}}$
& \cref{sec:11.7} \\
Single-phonon absorption (\cref{lim:11.8a}), \cref{eq:11.8-Gamma,eq:11.8-G,eq:11.8-Rscalar}
& \cite[Eqs.~(2), (7), (14), (16), (40)--(44)]{mitridate2024effective}; \cite[Eqs.~(2), (3)]{knapen2018detection}
& core kinetic row $\sim10^{-4}$ per u in $d$ (estimate); isotopic mass \SI{5}{\percent} in $d$; static vertex \SI{0.3}{\percent}
& \cref{sec:11.8} \\
Bulk acoustic resonator (\cref{lim:11.8b}), \cref{eq:11.8-Pn}
& \cite[Eqs.~(2), (3), (5)--(8)]{trickle2025piezoelectric}; the source writes $m_V^2$ for $m_\phi\omega_n$ (equal on resonance)
& $\times1.8$ at $\mathfrak{Q}_{\mathrm{p}} = \mathfrak{Q}_{\mathrm{DM}}$; $\le$\,\SI{10}{\percent} for $\mathfrak{Q}_{\mathrm{p}}/\mathfrak{Q}_{\mathrm{DM}} > 7.9$ or $< 0.2$
& \cref{sec:11.8} \\
Resonant-mass drive (\cref{lim:11.8c}), \cref{eq:11.8-h,eq:11.8-drive}
& \cite[Eqs.~(1), (3), (4), (6)--(8)]{arvanitaki2016sound}
& relativistic $\dd\ln(a/a_0)/\dd\ln\alpha$: $-1.97\times10^{-2}$ (Cu), $-2.2\times10^{-4}$ (Si)
& \cref{sec:11.8} \\
Atomic Migdal effect at dipole order (\cref{lim:12.2a}), \cref{eq:12.2-joint,eq:12.2-dipole,eq:12.2-shell,eq:12.2-rate}
& \cite[Eqs.~(55), (57), (73), (86)--(92)]{ibe2018migdal}; tables of \cite{cox2023precise}
& $O(\xi_D^2)$: $\le$\,\SI{10}{\percent} for $\xi_D \le 0.127$ (Ge), 0.188 (Si)
& \cref{sec:12.2} \\
Migdal excitation of bound final states (\cref{lim:12.2b}), \cref{eq:12.2-exc,eq:12.2-bound}
& \cite[Eqs.~(1), (6)]{kahlhoefer2026probing}; oscillator strengths \cite{kramida2024nist}
& $\le$\,\SI{10}{\percent} for $\xi_D \le 0.24$ (H); lines $n \le 3$ hold \SI{84}{\percent} (He)
& \cref{sec:12.2} \\
Crystal Migdal effect through the loss function, soft limit (\cref{lim:12.3a}), \cref{eq:12.3-soft,eq:12.3-LMig}
& \cite[Eq.~(6)]{knapen2021migdal}
& $\times1.35$ (local fields, Si); free-recoil bracket 0.04 (\SI{100}{eV}) to 0.42 (\SI{1}{keV})
& \cref{sec:12.3} \\
Migdal effect in semiconductors with local fields, free recoil (\cref{lim:12.3b}), \cref{eq:12.3-FL,eq:12.3-rate}
& \cite[Eqs.~(1), (7), (11), (12)]{liang2021describing}
& $1.1\times10^{-3}$ for $\mchi = 3$--\SI{150}{MeV}; 0.91 at \SI{1}{GeV}; probability 0.62 at $E_R = \SI{500}{eV}$
& \cref{sec:12.3} \\
Per-mass Migdal rate with the valence ionization factor (\cref{lim:12.3c}), \cref{eq:12.3-permass}
& \cite[Eqs.~(1), (6), (A36)]{knapen2021migdal}
& $\le$\,\SI{10}{\percent} for $\mchi \ge \SI{66}{MeV}$; $\times1.25$ at 30, $\times1.85$ at \SI{20}{MeV}
& \cref{sec:12.3} \\
Wannier picture of the boost (\cref{lim:12.3d}), \cref{eq:12.3-AW,eq:12.3-AWbound}
& \cite[Eqs.~(2.7), (2.8)]{liang2020describing}
& unsized
& \cref{sec:12.3} \\
Migdal effect with phonon final states (\cref{lim:12.4a}), \cref{eq:12.4-harm,eq:12.4-rate}
& \cite[Eqs.~(3.1), (3.6), (3.8), (4.1), (4.5)]{berghaus2023migdal}; \cite[Eq.~(21)]{esposito2025migdal}; \cite[Eqs.~(3.11), (4.1), (4.6)--(4.8)]{berghaus2026migdal}
& $\lesssim 3\times10^{-3}$ (Si, $\omega_e = \SI{1.2}{eV}$)
& \cref{sec:12.4} \\
Phonon-mediated Migdal effect (\cref{lim:12.4b}): \cref{lim:12.4a} with the virtual-phonon propagator, \cref{eq:12.4-pmrate}
& \cite[Eqs.~(3.1)--(3.4)]{liang2022phonon}
& $2\langle\omega_{\mathrm{ph}}^2\rangle/\omega_e^2 = 2.9\times10^{-3}$ at the Si gap; at order $p^2$ phonon-emitting fraction $\times2.4$ (\SI{5}{keV})
& \cref{sec:12.4} \\
Migdal-type absorption of a scalar coupled to nucleons (\cref{lim:12.6a}), \cref{eq:12.6-absrate}
& \cite[Eqs.~(13)--(17)]{dzuba2024migdal}
& $\tfrac32\omega_p^2/\omega^2$: $1.2\times10^{-3}$ at \SI{1}{keV}; $\le$\,\SI{10}{\percent} above \SI{57}{eV}
& \cref{sec:12.6} \\
Phonon-assisted electron-recoil rate and loss function (no box), \cref{eq:13.3-SPA,eq:13.3-LPA}; onset exponent, \cref{eq:13.3-pon}
& ---
& $\theta_{\mathrm{ep}} = 0.010$--$0.016$ (Si, absorption edge)
& \cref{sec:13.3} \\
Phonon-assisted absorption as the optical limit (\cref{lim:13.5a}), \cref{eq:13.5-limit}; edge \cref{eq:6.5-edge}
& \cite[Eq.~(1)]{noffsinger2012phonon}; \cite{hall1954infrared}
& data/edge $= 1.4$ at \SI{36}{meV}, 2 at \SI{17}{meV} above $\Eg$ (Si, \SI{300}{K})
& \cref{sec:13.5} \\
Inverse temperature transfer of absorption data (no box), \cref{eq:13.5-forward,eq:13.5-inverse}
& ---
& $\approx\SI{12}{\percent}$ from \SI{77}{K}, $\approx\SI{40}{\percent}$ from \SI{295}{K}
& \cref{sec:13.5} \\
Zero-point renormalization by special displacements (\cref{lim:13.7a}), \cref{eq:13.7-WL,eq:E.8-zg}; adiabatic AHC \cref{eq:6.3-adiabatic}
& \cite[Eqs.~(3)--(6), (10), (20), (24)]{zacharias2016one}; \cite{allen1981temperature}
& fourth order \SI{1}{\percent} ($-\SI{0.4}{meV}$, model); nonadiabatic \SI{8}{meV} on the Si gap; sideband $(1 + \omega_\nu/x)^3$ near the edge
& \cref{sec:13.7} \\
Nonadiabatic correction to the zero-point renormalization (no box), \cref{eq:13.7-na,eq:13.7-sqrt}
& ---
& factor $\approx2$ (estimate against calculation)
& \cref{sec:13.7} \\
One-magnon rate at linear spin-wave order (\cref{lim:14.1a}), \cref{eq:14.1-Safm,eq:14.1-rate}
& \cites[Eqs.~(5)--(12)]{trickle2020detecting}[Eqs.~(34), (38)]{trickle2022effective}
& about \SI{7.8}{\percent} for NiO; \SI{3.1}{\percent} at $|\vect{S}| = \frac52$, \SI{16}{\percent} at $|\vect{S}| = \frac12$
& \cref{sec:14.1} \\
Longitudinal and transverse spin weights (no box), \cref{eq:14.1-ratio}
& ---
& unsized
& \cref{sec:14.1} \\
Absorption into the uniform mode: the Larmor line (\cref{lim:14.3a}, panel a), \cref{eq:14.3-larmor,eq:14.3-rate}; $\lambda_S = g_J - 1$, \cref{eq:14.3-gJ}
& \cites[Eqs.~(2.4), (2.5)]{oshikawa2002electron}[Eqs.~(27), (30), (34), (35)]{mitridate2020detectability}[Eqs.~(25)--(27)]{chigusa2020detecting}
& $\le$\,\SI{10}{\percent} in line position and weight for $B_0 \ge 4.8\,\mu_0M_{\mathrm{s}}$ (thin film, field in plane); none for a saturated sphere
& \cref{sec:14.3} \\
Antiferromagnetic-resonance weight (\cref{lim:14.3a}, panel b), \cref{eq:14.3-afmr,eq:14.3-afmrweight}
& \cite[Eqs.~(2.4), (2.5)]{oshikawa2002electron} (Larmor statement, which this weight replaces)
& zero-point correction 5--\SI{23}{\percent} for $|\vect{S}| = \frac32$ at $\omega_A/\omega_E = 2$ to 0.3 (eight-site cluster); shape $\sim\omega_M/\omega_E = 1.6\times10^{-2}$ (MnF$_2$)
& \cref{sec:14.3} \\
Data-driven spin response (\cref{lim:14.4a}), \cref{eq:14.4-PT,eq:14.4-data}
& \cite[Eqs.~(8), (11), (26); dictionary Eq.~(B4)]{berlin2025determining}
& $-$\SI{17}{\percent} for Ni$^{2+}$ at $g_{\mathrm{ion}} = 2.2$ ($C_{\mathrm{sn}}(0) = 0.83$); $-$\SI{27}{\percent} at $g_{\mathrm{ion}} = 2.34$
& \cref{sec:14.4} \\
Chiral-phonon detection probability as the band average of $|\vect{L}^{\mathrm{ph}}|$ (\cref{lim:14.5a}, panel a), \cref{eq:14.5-Lph,eq:14.5-muph}
& \cites[Eqs.~(6), (7)]{romao2023chiral}[Eqs.~(3), (4), (11), (12)]{matas2025chiral}
& $-$\SI{10}{\percent} for $\mchi \ge \SI{2.2}{MeV}$ (GaAs, lowest band); below \SI{1}{\percent} of the rate in that band for $\mchi < \SI{0.8}{MeV}$; force constants 0.124 against 0.07, grid, masses and weighting $\le$\,\SI{14}{\percent}, long-range dipole part $-$\SI{14}{\percent} to $+$\SI{40}{\percent}
& \cref{sec:14.5} \\
Metastable line of a single-molecule magnet (\cref{lim:14.5a}, panel b), \cref{eq:14.5-smm,eq:14.5-meta}
& \cite[Eqs.~(3), (4); rate Eqs.~(9), (10), (17), (18)]{alves2026search}
& $-$\SI{8}{\percent} at \SI{4}{K}, $-$\SI{38}{\percent} at \SI{8}{K} (Mn$_{12}$-acetate, \SI{1}{T})
& \cref{sec:14.5} \\
Static magnetoelectric response of a multiferroic (\cref{lim:14.5b}), \cref{eq:14.5-cross,eq:14.5-GL}
& \cite[Eqs.~(5), (10)]{roising2021axion}
& Landau over spin waves 1.4 at $0.3\,T_{\mathrm{CM}}$, 4.5 at $0.1\,T_{\mathrm{CM}}$ (model); no size near $T_{\mathrm{CM}}$
& \cref{sec:14.5} \\
Transverse magnetization of a driven nuclear-spin ensemble (\cref{lim:14.6a}), \cref{eq:14.6-mperp}
& \cites[Eq.~(1)]{budker2014proposal}[Eq.~(1) and Eqs.~(S22)--(S25), (S41), (S53)]{aybas2021search}
& $\le$\,\SI{10}{\percent} for $t \le 1.6\,t_{\mathrm{coh}}$ at $T_2 = 10\,t_{\mathrm{coh}}$; \SI{0.6}{\percent} at all $t$ for $T_2 = 0.1\,t_{\mathrm{coh}}$
& \cref{sec:14.6} \\
Gaussian-field signal and its power spectrum (\cref{lim:14.6b}), \cref{eq:14.6-Sphi,eq:14.6-Fphi,eq:14.6-expbin}
& \cite[Eqs.~(8), (10), (24)]{foster2018revealing}
& $\le$\,\SI{10}{\percent} for $t_{\mathrm{int}} \ge 13\,t_{\mathrm{coh}}$ (sharp target); $1 - 1.25\,t_{\mathrm{coh}}/t_{\mathrm{int}}$ above
& \cref{sec:14.6} \\
Sensitivity scalings of the exponential-bin likelihood (\cref{lim:14.6c}), \cref{eq:14.6-like,eq:14.6-scaling}
& \cites[Eqs.~(A7), (A8)]{budker2014proposal}[Eqs.~(29), (54), (63), (69)]{foster2018revealing}
& fixed amplitude too strong by 1.99 in the coupling below $t_{\mathrm{coh}}$, 1.76 at $t_{\mathrm{coh}}$, 1.015 at $100\,t_{\mathrm{coh}}$
& \cref{sec:14.6} \\
Loss function of a superconductor: conserving response $\to$ bare bubble $\to$ $\mathfrak{g}_{\mathrm{I}}L^{\mathrm{nor}}$ $\to$ Lindhard (\cref{lim:15.1a}), \cref{eq:15.1-Lcons,eq:15.1-Rcons,eq:15.1-gI}
& \cites[Eqs.~(1), (S.13)]{hochberg2021determining}[Eq.~(3.4)]{hochberg2016detecting}[Eqs.~(8)--(10)]{gan2026mean}[Eq.~(2) and text]{hochberg2023directional}
& $\mathfrak{g}_{\mathrm{I}}(y) - 1$: 0.15 at $y = 5$, 0.05 at $y = 10$; ladder: $L^{\mathrm{cons}}/(\mathfrak{g}_{\mathrm{I}}L^{\mathrm{nor}}) \in [0.980, 1.085]$ for $y \ge 3$ (clean)
& \cref{sec:15.1} \\
First BCS correction to the screening (no box), \cref{eq:15.1-sBCS}
& ---
& bare bubble at the edge 0.935 vs 0.961 ($y = 10$), 0.796 vs 0.835 ($y = 5$)
& \cref{sec:15.1} \\
Resummed loss function on the kinematic edge (\cref{lim:15.2a}), \cref{eq:15.2-LBCS,eq:15.2-cons}
& \cite[Eqs.~(8), (9)]{gan2026mean}
& bare over conserving 0.38, 0.53, 0.75, 0.92 at $y = 2.2$, 3, 5, 10 on the edge
& \cref{sec:15.2} \\
Drude and Mermin loss functions of a metal with collisions (no box), \cref{eq:15.3-limits,eq:15.3-mermin}
& \cite[Eq.~(19)]{knapen2021dark}; \cite[Eqs.~(168), (172)]{kahn2022searches}; \cite[Eqs.~(S.17), (S.18)]{hochberg2021determining}; \cite{mermin1970lindhard}
& 1.43 against the ballistic 0.79 at $q\mathfrak{l}_e = 1.5$; \SI{0.2}{\percent} for the Drude form at $q\mathfrak{l}_e = 0.086$
& \cref{sec:15.3} \\
Phonon-assisted (Holstein) absorption of a superconductor (\cref{lim:15.4a}), \cref{eq:15.4-T,eq:15.4-holstein}
& \cite[Eqs.~(13), (16), (17)]{hochberg2016absorption}; Allen~\cite{allen1971electron}, as quoted there
& $\sigma_1^{\mathrm{S}}/\sigma_1^{\mathrm{N}} = 0.42$ ($\omega'^4$) or 0.73 ($\omega'^2$) at \SI{1}{meV}; $\times0.52$ from $1/(1 + \lambda_{\mathrm{tr}})^2$, known for the normal metal at $\omega \ll \omega_{\mathrm{D}}$~\cite[Eq.~(3)]{allen2015electron}, open for the superconducting form
& \cref{sec:15.4} \\
Impurity-assisted (Mattis--Bardeen) absorption (no box), \cref{eq:15.4-MB,eq:D.40-g}
& \cite{mattis1958theory}
& unsized
& \cref{sec:15.4} \\
Spin-current absorption in a clean superconductor (no box), \cref{eq:15.4-sJdef,eq:15.4-sJ}
& ---
& $g_{aee} = 5.4$--$7.1\times10^{-13}$ for $|F_1^a/3| \le 0.1$ (Leggett form $(1 + F)(1 - f)/(1 + Ff)$)
& \cref{sec:15.4} \\
Dirac-material rates with the background in the total $\varepsilon$ (\cref{lim:16.1a}), \cref{eq:16.1-L,eq:16.1-ratio}
& \cite[Eqs.~(2.6), (3.7), (3.11), (3.13), (3.14), (4.4), (4.5)]{hochberg2018detection}; \cite[Eqs.~(2.5)--(2.7), (2.25)--(2.27), (3.3), (3.13), (3.14)]{geilhufe2020dirac}
& published form $\times\varepsilon_{\mathrm{bg}}^2$ (scattering), $\times\varepsilon_{\mathrm{bg}}$ (absorption): 1600 and 40 at $\varepsilon_{\mathrm{bg}} = 40$; 5200--6400 for ZrTe$_5$ scattering (measured-structure row, $\Delta_{\mathrm{D}} = \SI{11.8}{meV}$, book halo); 4800--6600 over the rows at 11.8 and \SI{15.6}{meV} and the halos (238, 250.5, 544) and (220, 232, 600)~km\,s$^{-1}$; 4300 with $v_0 = 220$, $v_{\mathrm{esc}} = \SI{544}{km.s^{-1}}$
& \cref{sec:16.1} \\
One-dimensional plasmon of linear channels (\cref{lim:16.2a}), \cref{eq:16.2-wire,eq:16.2-kernel1D}
& \cite[Eqs.~(3)--(6)]{zhang2026luttinger}
& $\le$\,\SI{10}{\percent} for $q_\parallel < 1.36\,k_{\mathrm{F}}$; $+$\SI{4.0}{\percent} in $\omega_{\mathrm{pl}}$ at $q_\parallel = k_{\mathrm{F}}$
& \cref{sec:16.2} \\
Acoustic (demon) plasmon of two components and its Landau width (\cref{lim:16.2b}), \cref{eq:16.2-demon,eq:16.2-share,eq:16.2-width}
& \cite{pines1956electron}; \cite[Eqs.~(1), (2)]{hochberg2025unconventional}
& $\le$\,\SI{10}{\percent} for $v_a/v_{\mathrm{F},l} < 0.157$ (3D) and $< 0.53$ (2D); $+$\SI{3.5}{\percent} in $\omega_a$ from the heavy band at $v_a/v_{\mathrm{F},h} = 3.16$
& \cref{sec:16.2} \\
Effective medium of a layered stack (\cref{lim:16.2c}), \cref{eq:16.2-emt,eq:16.2-inplane,eq:16.2-corr}
& \cite[Eq.~(16); thin film Eq.~(19)]{lasenby2022dark}
& $\le$\,\SI{10}{\percent} for $ka_s < 1.94$ at \SI{30}{\degree} and $< 7.2$ in the plane; 0.7 and \SI{2.8}{\percent} at $ka_s = 1$
& \cref{sec:16.2} \\
Vacuum ejection from graphene and nanotubes with plane-wave final states (\cref{lim:16.3a}), \cref{eq:16.3-omega,eq:16.3-Sej,eq:16.3-rate,eq:16.3-iso}
& \cite[Eqs.~(1)--(5), (7)]{hochberg2017directional}; \cite[Eqs.~(16)--(18), (23)]{cavoto2018sub}
& OPW factor 1.1 to 3.5 for $0.6 < E_{\mathrm{ej}} < \SI{94}{eV}$ (4.3 at the end of the \SI{10}{MeV} spectrum); total rate $\times1.44$ (\SI{10}{MeV}) and $\times1.84$ (\SI{100}{MeV}), heavy mediator; final-state factor between 1 and 7.3, 3.4, 2.6 at 10, 50, \SI{100}{eV}
& \cref{sec:16.3} \\
Interband response of a semiconducting nanotube, one-dimensional massive Dirac bands in the random-phase approximation (no box), \cref{eq:16.3-R1D,eq:16.3-eps2}
& ---
& unsized
& \cref{sec:16.3} \\
Quantum-dot scattering rate and its bulk limit, with the prefactor of the published rate corrected (\cref{lim:16.3b}), \cref{eq:16.3-gap,eq:16.3-qdrate}
& \cite[Eqs.~(2)--(6), (10)--(12)]{blanco2023dark}; \cite[Eqs.~(3.13), (3.17)]{essig2016direct}
& 0.76 to 1.03 of the scissored bulk at 3--11 confinement energies ($qR = 8$; 0.70--0.97 at $qR = 3$), 2.4--66 nearer the gap (hard-wall model); within \SI{10}{\percent} only above $109\,E_{\mathrm{conf}}$
& \cref{sec:16.3} \\
Dark-photon absorption by quantum dots in a host: the depolarization factor (no box), \cref{eq:16.3-depol}
& \cite[Eq.~(16)]{blanco2023dark}
& $\times6.00$ for PbS ($\varepsilon_{\mathrm{in}} = 17.63$) in silica; 4.62--6.71 for $\Real\varepsilon_{\mathrm{in}} = 10$--25; 0.566 for an illustrative $\varepsilon_{\mathrm{in}} = 1 + \ii$ at \SI{5}{eV}
& \cref{sec:16.3} \\
Frenkel-exciton rate binned in $\omega$: incoherent sum of molecular form factors (\cref{lim:16.4a}), \cref{eq:16.4-ff,eq:16.4-binned,eq:16.4-rate}
& \cites[Eqs.~(11), (13), (16)--(19)]{blanco2020dark}[Eqs.~(8), (9)]{blanco2021dark}
& $+6.3$\,\% (heavy, $\mchi \ge \SI{30}{MeV}$), up to $+27$\,\% between 1.2 and \SI{7.5}{MeV}, $+38$\,\% at 0.88--\SI{1.0}{MeV}, $+29$ to $+35$\,\% (light, $\mchi \ge \SI{2}{MeV}$), in a $\pi$-electron model
& \cref{sec:16.4} \\
Dark-photon absorption by dilute donors, host screening counted once (\cref{lim:16.5a}), \cref{eq:16.5-SD,eq:16.5-abs,eq:16.5-gos}
& \cite[Eqs.~(9), (10), (12) and the absorption rate after Eq.~(12)]{du2022doped}
& $\le$\,\SI{10}{\percent} for $\nd_{\mathrm{D}} < \SI{2.6e17}{cm^{-3}}$ (continuum, both ground states; \SI{8.8e17}{cm^{-3}} with the central cell); at a \SI{4.2}{\micro eV} line center only below \SI{4.8e15}{cm^{-3}}; impurity band at $\nd_{\mathrm{D}}/\nd_{\mathrm{c}} \approx 0.3$--0.6 unsized
& \cref{sec:16.5} \\
Color-center creation by a nuclear recoil, free recoil with a step at the displacement threshold (\cref{lim:16.6a}, panel a), \cref{eq:16.6-yield,eq:16.6-ccrate}
& \cite[Eqs.~(1), (2), (5)--(7)]{budnik2018direct}
& rate ratio within \SI{10}{\percent} for $x \ge 1.36$ (\SI{77.2}{meV}) and 1.22 (\SI{33.5}{meV}), within \SI{2}{\percent} for $x \ge 1.96$; nonzero below $x = 1$; directional threshold ($a = 0.4$) away from the step by more than \SI{10}{\percent} up to $x = 4.8$
& \cref{sec:16.6} \\
Bond breaking in a molecular solid, classical free-recoil rate with the dissociation threshold (\cref{lim:16.6a}, panel b), \cref{eq:16.6-conv,eq:16.6-qmin,eq:16.6-diss}
& \cite[Eqs.~(3.3), (3.11), (3.12), (3.17), (4.3), (4.5)]{essig2017detection}
& within \SI{10}{\percent} for $x \ge 2.65$ (impulse) and 3.61 (improved Born, $\mchi \ge \SI{50.5}{MeV}$); mass scale lowered by a factor 0.84 (impulse) and 0.75 (Born) at the classical rate of $x = 1.1$
& \cref{sec:16.6} \\
Two-level response to a classical strain or field, $\tan\delta = \tan\delta_0\tanh(\omega/2T)$, $Q^{-1} = \pi C_j\tanh(\omega/2T)$ (\cref{lim:16.7a}), \cref{eq:16.7-chires,eq:16.7-tand}
& \cites[Eqs.~(4), (12), (13), (15)]{meissner2021acoustic}[Eq.~(20), electric form]{muller2019towards}
& $\le$\,\SI{10}{\percent} for $T < \SI{1.36}{K}$ at \SI{5}{GHz} and $< \SI{0.43}{K}$ at \SI{0.5}{GHz}; $4\times10^{-9}$ at \SI{10}{mK}, \SI{5}{GHz}
& \cref{sec:16.7} \\
Dark-photon absorption on two-level systems (\cref{lim:16.7b}), \cref{eq:16.7-dp}
& \cite[Eq.~(20), electric form]{muller2019towards}
& $\le$\,\SI{10}{\percent} for $m_\phi > \SI{1.15}{neV}$ at \SI{10}{mK} and $> \SI{115}{neV}$ at \SI{100}{mK}; $\tan^2\delta_0 \approx 10^{-7}$
& \cref{sec:16.7} \\
Scattering off a random medium in the first Born approximation and its validity (\cref{lim:16.7c}), \cref{eq:16.7-born,eq:16.7-Ffluc,eq:16.7-scrit}
& \cite[Eqs.~(2), (9)--(12), (14)]{liu2025coherence}
& $\le$\,\SI{10}{\percent} for $|V_{\mathrm{opt}}|/E_n < 0.092$ (repulsive), $< 0.110$ (attractive); relative correction $\pm V_{\mathrm{opt}}/E_n$
& \cref{sec:16.7} \\
Glass absorption below the Ioffe--Regel frequency, $C^{\mathrm{IR}} = A^{\mathrm{IR}} + B^{\mathrm{IR}}\omega^2$ (\cref{lim:16.7d}), \cref{eq:16.7-eps2,eq:16.7-CIR,eq:16.7-AB}
& \cites[Eq.~(1)]{taraskin2006universal}[Eq.~(2)]{bloch2026dark}
& $B^{\mathrm{IR}}\omega^2/A^{\mathrm{IR}} \le \SI{9}{\percent}$ for $\omega < \omega_{\mathrm{IoR}}$ in silica
& \cref{sec:16.7} \\
Two-level relaxation by one phonon (Jäckle) and by conduction electrons (Black), with the superconducting factor (\cref{lim:16.8a}, panel a), \cref{eq:16.8-golden,eq:16.8-Jph,eq:16.7-tau1,eq:16.8-black,eq:16.8-sc}
& \cites[Eqs.~(7)--(9)]{meissner2021acoustic}[Eqs.~(4), (5)]{sahling2010anomalous}
& superconducting factor within \SI{10}{\percent} for $\Delta_{\mathrm{tl}} \le 0.69$--$0.73\,T$; 0.86 at $\Delta_{\mathrm{tl}} = T$, 0.71 at $2T$
& \cref{sec:16.8} \\
Ohmic coupling of a tunneling defect in the Born limit, $\alpha_{\mathrm{ohm}} = 4K_{\mathrm{b}}$ (\cref{lim:16.8a}, panel b), \cref{eq:16.8-born}
& \cite[Eq.~(4) and the single-spin density of states of its supplement]{maillet2023nanomechanical}
& $\le$\,\SI{10}{\percent} for $K_{\mathrm{b}} \le 0.017$; ratio 0.70 at $K_{\mathrm{b}} = 0.07$ and 0.60 at 0.11
& \cref{sec:16.8} \\
Ohmic scaling limit of a symmetric tunneling defect, small-coupling relaxation rate $\tilde\Delta_0^2/\gamma_0$ (\cref{lim:16.8b}), \cref{eq:16.8-scaling,eq:16.8-qureshi,eq:16.8-Tcoh}
& \cite[Eqs.~(8)--(10)]{qureshi1995dynamics}
& $\le$\,\SI{10}{\percent} for $\alpha_{\mathrm{ohm}} \le 0.091$; 0.71 at 0.28, 0.56 at 0.44
& \cref{sec:16.8} \\
Thermally activated hopping and tunneling through one barrier: transition-state theory with Wigner and parabolic-barrier transmission, instanton below the crossover (\cref{lim:16.8c}), \cref{eq:16.8-kappa,eq:16.8-Tx,eq:16.8-Tfrz}
& \cites[Eqs.~(1)--(4), (69)--(71)]{lawrence2024semiclassical}[Eq.~(8)]{sahling2010anomalous}
& Wigner within \SI{10}{\percent} for $T \ge 2.1\,T_\times$; exact/classical 6.2 at $T_\times$, 2.2 at $1.5\,T_\times$; instanton \SI{18}{\percent} low at 0.25--$0.5\,T_\times$
& \cref{sec:16.8} \\
Heat release of a glass cooled from $T_{\mathrm{hi}}$ to $T_{\mathrm{lo}}$ (\cref{lim:16.8d}), \cref{eq:16.8-heat}
& \cites[Eqs.~(1), (6)--(8)]{sahling2010anomalous}[Table~1]{sahling2014anomalous}
& $\le$\,\SI{10}{\percent} for $1.8\,\tau_{\min} \le t \le 7.1\times10^{10}\,\tau_{\min}$ (10.6 of 12.2 decades); within $2\times10^{-4}$ from $10^2$ to $10^8\,\tau_{\min}$
& \cref{sec:16.8} \\
Partial (valence) f-sum of the loss function (\cref{lim:17.5a}), \cref{eq:17.5-partial}
& \cite[Table~I]{ramanathan2020ionization} (measured $\omega_p^{\mathrm{val}}$)
& $-$\SI{3}{\percent} core screening, $+$\SI{3}{\percent} shell transfer (Si)
& \cref{sec:17.5} \\
Pair energy and Fano factor as cascade moments (\cref{lim:18.2a}), \cref{eq:18.2-moments,eq:18.2-decomp}
& \cite[Eqs.~(2)--(3)]{ramanathan2020ionization}
& $b_N\varepsilon_{\mathrm{eh}}/\omega$: \SI{10}{\percent} at $\omega = 10\Eg$; residual after it $\le$\,\SI{0.9}{\percent} there
& \cref{sec:18.2} \\
Linear mean yield (no box), \cref{eq:18.2-linear,eq:18.2-tagged}
& \cite[Chap.~XI]{feller1971introduction} (key renewal theorem)
& unsized
& \cref{sec:18.2} \\
Random-parking (R\'enyi) limit (\cref{lim:18.2b}), $\rho_\infty$ and $F_{\mathrm{park}}$
& \cite[Eqs.~(2), (19), (23), (44)]{subashiev2010correlation}
& $\pm0.8\,\lambda_{\mathrm{AB}}$ for $\lambda_{\mathrm{AB}} \le 0.003$; $-$\SI{3.2}{\percent} ($\rho$), $+$\SI{3.7}{\percent} ($F$) at 0.1
& \cref{sec:18.2} \\
Renewal (uncorrelated final-state) Fano estimate (\cref{lim:18.2c}), \cref{eq:18.2-ren,eq:18.2-rencasc}
& \cite[Eq.~(7)]{subashiev2010correlation}
& $F_{\mathrm{eh}}/F_{\mathrm{ren}} = 1.15$ (parking) to 5.2 (Si), 7.9 (Ge)
& \cref{sec:18.2} \\
Step model and Ramanathan--Kurinsky model as limits of the seeded cascade distribution (\cref{lim:18.3a}), \cref{eq:18.3-step,eq:18.3-beta,eq:18.3-pq}
& \cites[Eqs.~(5)--(10), (14)]{ramanathan2020ionization}[Eq.~(5.1)]{essig2016direct}
& $d_{\mathrm{TV}} \le 0.062$ (energy-only model, one $\vect{k}$ grid; unresolved in the full band, where the photon seed moves 0.11--0.14 between grids) for a photon seed in place of the dark-matter seed (largest at \SI{5.67}{eV}; mean \SI{4.6}{\percent}); splits with $\alpha_{\mathrm{RK}} = 1$, $\to\infty$ up to 0.21, 0.39, one-sided split 0.67; step model up to 0.86
& \cref{sec:18.3} \\
Klein relation (\cref{lim:18.4a}), \cref{eq:18.4-shockley,eq:18.4-uniform,eq:18.4-klein}
& \cite[Eq.~(3)]{wei2016average}; \cite[Eq.~(1), Sect.~II]{nielsen2026first}
& $+$\SI{36}{\percent} at $\lambda_{\mathrm{AB}} \to 0$ ($\theta_{\mathrm{ii}} = \frac32$), $+$\SI{39}{\percent} at $\theta_{\mathrm{ii}} = 1$
& \cref{sec:18.4} \\
Monte Carlo planes of the random-$k$ cascade (\cref{lim:18.4b}), \cref{eq:18.4-RKplanes}
& \cite[Eqs.~(8), (12)--(13)]{ramanathan2020ionization}
& $-$1.3 to $-$\SI{2.4}{\percent} in $F$ at the five points; $-$6 to $-$\SI{7}{\percent} at $A_{\mathrm{AB}} = 1$, $-$\SI{4}{\percent} at 12
& \cref{sec:18.4} \\
Linear law of the pair energy (\cref{lim:18.4c}), \cref{eq:18.4-fit,eq:18.4-linear}
& \cite[Sect.~II]{nielsen2026first} (empirical $1.54\Eg + \SI{1.89}{eV}$)
& residuals $-19.8$ to $+\SI{7.0}{\percent}$ (rms \SI{10.5}{\percent}) against measurement
& \cref{sec:18.4} \\
Binary electron--hole threshold of a slow recoil (no box), \cref{eq:19.2-transfer,eq:19.2-Eth}
& ---
& $-$\SI{4.4}{\percent} (Si), $-$\SI{3.1}{\percent} (Ge) from the $v_N^2$ term; momentum tail unsized
& \cref{sec:19.2} \\
Ionization yield of silicon recoils below the binary threshold (no box), \cref{eq:19.2-Yq,eq:19.2-charge,eq:19.2-FNR}
& \cite[Table~II, Eq.~(1)]{albakry2023first}
& unsized
& \cref{sec:19.2} \\
Lindhard quenching (\cref{lim:19.3a}), \cref{eq:16.6-lindhard,eq:19.3-simple}
& \cite[Eqs.~(3)--(7)]{sarkis2020study}; \cite[Eqs.~(5)--(7)]{sorensen2015atomic}; \cite[Eq.~(3)]{agnese2018nuclear}; \cite{lindhard1963integral,lewin1996review}
& $-$8 to $-$\SI{38}{\percent} (data over formula, 0.1--\SI{3.9}{keV}); fit within 3--\SI{5}{\percent} of its own equation
& \cref{sec:19.3} \\
Bethe stopping from the loss function (\cref{lim:19.3b}), \cref{eq:19.3-bethe,eq:19.3-lnI}
& \cite[Eq.~(34.5), Sect.~34.2.1]{navas2024review}
& $\le$\SI{10}{\percent} for $v \ge 1.84v_{\mathrm{F}}$
& \cref{sec:19.3} \\
Displacement count and defect energy loss (\cref{lim:19.4a}), \cref{eq:19.4-count,eq:19.4-KP,eq:19.4-loss}
& \cite[Eq.~(1)]{nordlund2018improving}
& $+$11 to $+$\SI{14}{\percent} (formula over screened count, $T_{\mathrm{d}} \ge 10E_{\mathrm{d}}$); $-$\SI{20}{\percent} against hard spheres
& \cref{sec:19.4} \\
Unit heat quenching (\cref{lim:19.5a}), $Y_{\mathrm{ph}}(0) = 1$; \cref{eq:19.5-heatq,eq:19.5-release}
& \cite[p.~4]{alkhatib2021light}
& $+$\SI{41}{\percent} at $f_{\mathrm{ret}} = 0$ (\SI{100}{eV}); $+$\SI{1.7}{\percent} for \SI{5.9}{\percent} of the gap energy trapped or radiated
& \cref{sec:19.5} \\
Stored energy from the phonon-yield ratio (\cref{lim:19.5b}), \cref{eq:19.5-identity,eq:19.5-first,eq:19.5-bound,eq:19.5-widen}
& \cite[Eq.~(4), Table~II]{agnese2018energy}; \cite[Eqs.~(1), (4), Sects.~2, 3.2, 4]{agnese2018nuclear}
& $-0.0064$ in $Y_{\mathrm{ph}}$ from $\varepsilon_{\mathrm{cal}} \ne \varepsilon_{\mathrm{eh}}$ (Si, \SI{4}{V}); bound 0.081 at $\mathfrak{r}_{\mathrm{coll}} = 0.9$; first order against exact $\le 3.0\times10^{-3}$ for deficits in $[0, 0.05]$
& \cref{sec:19.4,sec:19.5} \\
Gap-energy release law (no box), \cref{eq:19.5-release,eq:19.5-bimol}
& ---
& unsized
& \cref{sec:19.5} \\
Two-bias inversion (no box), \cref{eq:19.5-linear,eq:D.45-inversion}
& ---
& unsized
& \cref{sec:19.5} \\
Defecton diffusion in a phonon gas (\cref{lim:19.6a}), \cref{eq:19.6-q2,eq:19.6-D}
& \cite[Eqs.~(12), (49)--(61), (64)]{pushkarov2004quantum}
& $\le$\SI{10}{\percent} for $T \ge 0.28\,\Delta_{\mathrm{dfn}}$
& \cref{sec:19.6} \\
Crowdion profile and kink mass (\cref{lim:19.6b}), \cref{eq:19.6-Mcr}; \cref{eq:C.20-kink,eq:C.20-Ekink,eq:C.20-EPN}
& \cite[Eqs.~(2), (3)]{fitzgerald2018structure}; \cite{braun2004frenkel}
& $-$\SI{1.4}{\percent} in $E_{\mathrm{kink}}$ at $d_{\mathrm{kink}} = a_{\mathrm{ch}}$; $E_{\mathrm{PN}}$ $-$\SI{22}{\percent} there, $-$\SI{7.5}{\percent} at $2a_{\mathrm{ch}}$
& \cref{sec:19.6} \\
Hecht collection (\cref{lim:20.1a}), \cref{eq:20.1-hecht}; \cref{eq:20.1-hecht2,eq:20.1-qavg}
& \cite{hecht1932mechanismus}, \cite[Eqs.~(1)--(4)]{lachish2017electron}
& $-\mu_{\mathrm{mob}}\mathcal{E}/v_{\mathrm{sat}}$ (relative) at low bias; $-$\SI{8.9}{\percent} at $\mu_{\mathrm{mob}}\mathcal{E} = 0.1\,v_{\mathrm{sat}}$
& \cref{sec:20.1} \\
Bias work and the Luke phonon energy (\cref{lim:20.1b}), \cref{eq:20.1-work}; \cref{eq:19.5-Wbias}
& \cite[Eq.~(1)]{agnese2018nuclear}
& $-d_{\mathrm{det}}/6\lambda_{\mathrm{d}}^e - d_{\mathrm{det}}/6\lambda_{\mathrm{d}}^h$; $-$\SI{3.3}{\percent} at $d_{\mathrm{det}} = 0.1\,\lambda_{\mathrm{d}}$ (equal drift lengths)
& \cref{sec:20.1} \\
Fano factor of the measured charge (\cref{lim:20.1c}), \cref{eq:20.1-Fpos}; \cref{eq:C.11-thinfano}
& \cite[Eqs.~(A6)--(A7)]{ramanathan2020ionization}
& $+F_{\mathrm{eh}}(1/f_{\mathrm{det}} - f_{\mathrm{det}})$ against thinning, 0.012 at $f_{\mathrm{det}} = 0.95$
& \cref{sec:20.1} \\
Down-conversion and the scaling spectrum (\cref{lim:20.2a}), \cref{eq:20.2-anh,eq:20.2-Estar,eq:20.2-scaling}; \cref{eq:20.2-frag}
& \cite[Eqs.~(6), (8), Table~I]{leman2012invited}; \cite[Sect.~2.2.2]{kelsey2023g4cmp}; \cite[Tables~3, 4]{hernandez2025modeling}
& $\le$\SI{10}{\percent} for $t \ge \SI{0.46}{ns}$ (Si)
& \cref{sec:20.2} \\
Rayleigh law of isotope scattering (\cref{lim:20.2b}), \cref{eq:20.2-rayleigh}; \cref{eq:20.2-isoborn,eq:E.12-iso}
& \cite[Eq.~(4), Table~4]{hernandez2025modeling}
& $+9.6\,(E_{\mathrm{ath}}/\omega_{\mathrm{D}})^2$; $+$\SI{5}{\percent} at \SI{1}{THz}, $+$\SI{10}{\percent} at \SI{5.6}{meV}
& \cref{sec:20.2} \\
Transfer-kernel bound on observable differences (no box), \cref{eq:20.2-bound}
& ---
& bound 0.15 against the largest model difference 0.012 at \SI{1}{\micro s} (Si, the births of \cref{fig:20-spectrum}); it reduces to $f_{\mathrm{s}} = 0.017$ only if $E_{\mathrm{mem}}$ lies below every hard birth
& \cref{sec:20.2} \\
Impedance law as the averaged interface transmission (\cref{lim:20.3a}), \cref{eq:C.18-transmission}; \cref{eq:20.3-avg,eq:20.3-recip}
& \cite{swartz1989thermal}
& factor $\mathfrak{n}_{\mathrm{Si}}/\mathfrak{n}_{\mathrm{Al}} = 0.392$ on the escape from aluminum; the corrected form within \SI{6.9}{\percent} of the average over the plotted range
& \cref{sec:20.3} \\
Lumped film absorption (no box), \cref{eq:20.3-pabs,eq:20.3-tl}
& \cite[Eq.~(17)]{kelsey2023g4cmp}
& factor $1/\bar t_{\mathrm{esc}} = 2.72$ in the thin-film limit for aluminum on silicon, 1.82 for \SI{60}{nm} at $\mathfrak{l}_{\mathrm{f}} = v_st_0$
& \cref{sec:20.3} \\
Pair-breaking efficiency (\cref{lim:20.4a}), \cref{eq:20.4-below}; \cref{eq:20.4-split,eq:20.4-tpb,eq:20.4-emit,eq:20.4-ppb,eq:20.4-cascade}
& \cite[pp.~2, 5, Fig.~4]{guruswamy2014quasiparticle}; \cite{kaplan1976quasiparticle}
& $\le 0.011$ for $t_l/t_0 \le 4$; 0.036--0.041 at $t_l/t_0 = 8$ (kinetic equations above the cascade, 5--$10\Delta_{\mathrm{BCS}}$)
& \cref{sec:20.4} \\
Collection ratio of Luke phonons, one film type (no box), \cref{eq:20.5-sign}; mean-field kernels \cref{eq:20.5-Gabs,eq:20.5-meanfield}
& ---
& 0.008 from the extra bare-face loss of a two-film detector; $10^{-3}$ for a thick film
& \cref{sec:20.5} \\
High-voltage readout ceiling (no box), \cref{eq:20.5-hv}
& ---
& factor 1.23 between the binomial and mean-charge limits at \SI{220}{eV} (1.22 against 1.50)
& \cref{sec:20.5} \\
Mean inverse speed of the standard halo (\cref{lim:21.1a}), \cref{eq:D.1-eta1,eq:D.1-eta2}
& \cite[Eqs.~(B.4), (B.5)]{essig2016direct}
& $+\SI{1}{\percent}$ at \SI{430}{km.s^{-1}}, $+\SI{10}{\percent}$ at \SI{736}{km.s^{-1}}, $+\SI{20}{\percent}$ at \SI{760}{km.s^{-1}}
& \cref{sec:21.1} \\
Halo-independent bound for a speed-only response (\cref{lim:21.2a}), \cref{eq:21.2-hib}; \cref{eq:21.2-column,eq:21.2-multipole}
& \cite[Eqs.~(A1)--(A2)]{fox2011integrating}; \cite[Sect.~2]{chen2021halo}; \cite[Eqs.~(2.17), (2.18)]{bernreuther2024extracting}
& $+$2 to $+$\SI{15}{\percent} below \SI{560}{km.s^{-1}} (sampled from \SI{340}{km.s^{-1}}), at most $+\SI{2.0}{\percent}$ above; $+\SI{0.4}{\percent}$ at \SI{790}{km.s^{-1}}
& \cref{sec:21.2} \\
\end{longtable}
}

\printbibliography[heading=subbibliography]
